\section{Calibration}
\label{sec:calibration}

Validity has already been settled by the coordinates.  Calibration may
therefore concentrate on information: selecting, among admissible laws, one
that represents the quoted strip without inventing unsupported detail.

\subsection{A one-expiry objective}

Suppose a slice contains quotes $(k_i,\sigma_i)$ with
$w_i=\sigma_i^2\tau$.  The reference fit minimizes
\begin{equation}
\sum_i\left(
\frac{c(k_i;\theta)-\Black(k_i,w_i)}
     {\partial_\sigma\Black(k_i,w_i)+\eta}
\right)^2
+\alpha\sum_{n\ge4}n^{2\nu}a_n^2.
\label{eq:objective}
\end{equation}
Each element of \eqref{eq:objective} answers one practical problem.  The
residual is computed in price, so implied-volatility inversion stays outside
the optimization loop, where it cannot fail on an exploratory trial point.
Division by the Black vega $\partial_\sigma\Black=\phin(d_+)\sqrt\tau$ makes
a central residual a volatility error --- the unit in which quotes are
judged --- while the floor $\eta$ lets the metric degrade to a scaled price
error where vega vanishes, instead of assigning unbounded weight to a
negligible option.  The $n^{2\nu}$ ridge makes fine modes progressively more
expensive but leaves $a_2$ and $a_3$, the main curvature and asymmetry
modes, unpenalized.

The reference objective also contains a smooth row
$\log(1+e^{50(\lambda_+-0.90)})$.  The admissibility chart already prevents
the wall from being reached; this row is instead an explicit prior against a
nearly exhausted right-tail moment budget.

A quoted option is an interval, not necessarily a point.  Three targets are
useful.  A \emph{mid} fit uses \eqref{eq:objective} directly.  A \emph{band}
fit uses the outside-band hinge together with a weak mid anchor.  If
$[p_i^{\rm lo},p_i^{\rm hi}]$ is the price band and $\phi_i$ the vega
normalizer, its two residuals are
\begin{equation}
\frac{(c_i-p_i^{\rm hi})^++(p_i^{\rm lo}-c_i)^+}{\phi_i},
\qquad
\sqrt{0.05}\frac{c_i-p_i^{\rm mid}}{\phi_i}.
\label{eq:band}
\end{equation}
The squared hinge is continuously differentiable at a band edge, and the
weak mid row --- its weight deliberately small --- selects one point when a
wide band leaves a flat valley of equally admissible laws.
A \emph{haircut} fit first shrinks the volatility band by $h$:
\begin{equation}
\sigma_i^{\rm lo}=\min(\sigma_i^{\rm bid}+h,\sigma_i^{\rm mid}),
\qquad
\sigma_i^{\rm hi}=\max(\sigma_i^{\rm mid},\sigma_i^{\rm ask}-h).
\label{eq:haircut}
\end{equation}
If the original spread is narrower than $2h$, the band collapses to the mid
instead of changing orientation.  This occurs for every SPY quote in the
frozen sample; the wider NVDA spreads retain live bands.

The two underliers illustrate both regimes.  The SPY December node has median
spread \MacSpyDecMedianSpreadBp\ vol bp, far inside the default
$2h=100$ vol-bp width, and none of its haircut bands survives.  Its figures
therefore show raw bid--ask intervals around a mid fit.  On the featured NVDA
nodes, \MacNvdaOneDayBandLivePct\% and \MacNvdaLongBandLivePct\% of the
bands remain live.  The same rule becomes a mid target on a liquid index and a
genuine interval target on the wider single-name book.

The informational division can be read by location.  Between liquid quotes a
small change in the law moves observed prices and is strongly penalized: the
strip determines the call curve there, and with it the ATM level, skew,
curvature, and low-order moments.  Near the edge of the strip, data and
regularization share control; beyond the outer quote, the endpoint speeds
and the ridge supply most of the continuation.  The curve itself is smooth
--- what changes gradually is the source of information, a transition made
visible by plotting quote locations with the density as in
\cref{fig:lqdchart}.  Tail coordinates, ridge, and basis order should be
reported with a fit because they make these choices.

\subsection{Order and starting values}
\label{sec:orderguard}

At order $N$ the model has $N+1$ parameters.  The reference implementation
uses the effective order
\begin{equation}
N_{\rm eff}=\max\left(
\min\left(N,\left\lfloor\frac{n_q}{2}\right\rfloor-1\right),
\min(N,6)\right),
\label{eq:orderguard}
\end{equation}
where $n_q$ is the quote count.  The cap keeps the parameter count well below
the number of observations; the floor leaves thin strips enough shape to form
a usable smile.  The rule is an identification convention, not a validity
condition.  Changing it changes inter-quote density detail while every
resulting slice remains arbitrage-free.

Away from the floor, \eqref{eq:orderguard} requires roughly two quotes per
body degree of freedom.  This is not a theorem about optimal smoothing; it is
a guard against inferring one oscillatory direction from each noisy
observation.  On the shortest NVDA slice,
\MacGuardQuoteCountZeroDte\ quotes give
$N_{\rm eff}=\MacGuardEffOrderZeroDte$.  The fit remains visibly less certain
there rather than borrowing high-order detail from the basis.

Cold starts use the constant-speed value \eqref{eq:coldstart}.  In a maturity
sweep, the previous expiry supplies a warm start.  A multi-start check on the
featured SPY and NVDA nodes used \MacMultiStartCount\ randomized starts per
node and found \MacMultiStartBasins\ basin, with worst fit-quality spread
\MacMultiStartWorstDIvBp\ vol bp.  This supports the starting rule on those
slices; event-dominated books may still require a broader search.

\subsection{Coordinates for optimization}
\label{sec:charts}

Three charts describe the same polynomial family.

\begin{enumerate}
\item The \emph{raw chart} is $(L,R,a_2,\ldots,a_N)$.  It is simple, but
\eqref{eq:endpoints} couples body modes to both tails.
\item The \emph{endpoint chart} is
$(\log\lambda_-,\log\lambda_+,a_2,\ldots,a_N)$.  Solving
\eqref{eq:endpoints} for $(L,R)$ is linear.  Equivalently, each body basis
function is replaced by an endpoint-vanishing combination.  Body moves then
hold both tail scales fixed.
\item The \emph{logistic chart}, the reference default, writes
$\lambda_+=\expit\xi$ with $\xi\in\R$.  Every finite coordinate vector is
admissible, and
\begin{equation}
\frac{\dd\log\lambda_+}{\dd\xi}=1-\lambda_+.
\label{eq:chartcompression}
\end{equation}
Steps shrink smoothly as the wall is approached.
\end{enumerate}

The endpoint-vanishing body functions are explicit:
\begin{equation}
\widetilde P_n(u)=P_n(1-2u)-(1-u)-(-1)^nu,\qquad n\ge2.
\label{eq:endpointbasis}
\end{equation}
They satisfy $\widetilde P_n(0)=\widetilde P_n(1)=0$.  A tail policy can
therefore set $(\lambda_-,\lambda_+)$ while calibration adjusts the body in a
subspace that cannot move either endpoint.

The objective is identical in all three charts.  Fits in the frozen sample
agree across charts to \MacChartEquivParamAgree\ in the recovered raw
parameters.  In floating point, $\expit\xi$ eventually rounds to one, so the
implementation also enforces $\lambda_+\le1-10^{-6}$.

\subsection{Calibration sequence}

The complete one-slice calculation is short enough to state as an algorithm.

\begin{enumerate}
\item Convert strikes and prices to $(k_i,c_i)$ using the node's forward and
discount factor.  Map bid and ask volatilities through Black once to construct
price bands and compute the quote vegas used for normalization.
\item Choose the effective order from \eqref{eq:orderguard} and select mid,
band, or haircut residuals.  This fixes the data part of the objective before
optimization begins.
\item Construct a cold start from the ATM variance using
\eqref{eq:coldstart}, or map the previous expiry into the endpoint chart for a
warm start.  In either case, convert to logistic coordinates so all finite
trial vectors are admissible.
\item At each trust-region step, build $(x,x',G,G')$ on the fixed grid of
\cref{sec:numerics}, price every quote by \eqref{eq:call}, and obtain the
whole residual Jacobian from \eqref{eq:pass1}--\eqref{eq:pass3}.  Failed
exponent-budget trials return a large residual rather than an undefined
price.
\item After convergence, rebuild the accepted parameters on the dense grid.
Evaluate ATM handles, tail moments, Lee slopes, the log-contract variance, and
the numerical certificates (\cref{sec:certificates}) from this one final
object.
\item If the slice belongs to a maturity sequence, add calendar rows against
the accepted nearer slice and repeat.  Report any residual common-support
gap in normalized price units.
\end{enumerate}

This order matters.  Market conventions are fixed before the solve;
probability validity holds during the solve; diagnostics are computed after
the solve from the same law that produced the quoted prices.  No separate
post-processing smoother is inserted between calibration and publication.
